\section{Performance benchmarking}
Benchmarking is the comparison of one or more algorithms or queries across multiple metrics. Benchmarking enables fair, reproducible, and accurate comparisons in controlled environments with different features, configurations, and inputs \parencite{beyer_2015_benchmarking_and_resource_measurement, halkos2018_corporate_social_responsibility}. To develop a benchmarking framework that enables fair, reproducible, and accurate comparisons, several requirements must be met. It should ensure that data are obtained through valid and reliable measurements to ensure reproducible results. 

Results are reproducible if the same results can be obtained later by re-running the benchmarks on a system with the same configuration and specification as the previous benchmarks. To achieve this trait within a benchmarking framework, measurements must be reliable. A reliable measurement is both accurate and precise \parencite{beyer_2015_benchmarking_and_resource_measurement}. According to \textcite[pp.~21--22]{jcgm2012_international_vocabulary_of_metrology}, measurement accuracy is defined as ``closeness of agreement between a measured quantity value and a true quantity value of a measurand,''  while measurement precision is defined as ``closeness of agreement between indications or measured quantity values obtained by replicate measurement on the same or similar objects under specified conditions.'' 

% \subsection{Measuring run-time using CPU-time}
% A simple time measurement, such as elapsed time computed by a wall clock timer, is insufficient for short-running queries. These types of timers do not allow for comparing the resource usage of multi-threaded tools and may be unknowingly influenced by \acrshort{io}-operations. \acrshort{cpu}-time is a more meaningful measurement method and accounts for the aforementioned. However, using \acrshort{cpu}-time makes the time measurement more complex and introduces challenges, especially if child processes are involved \parencite{beyer_2015_benchmarking_and_resource_measurement}.

% The \acrshort{cpu}-time is a meaningful way to measure compared to wall time. Wall time \parencite{beyer_2015_benchmarking_and_resource_measurement}.

\subsection{Validity, reliability, and reproducibility}
\subsection{Variations and error sources in elapsed-time measurements}
 A common issue in run-time measurements is fluctuation in consecutive measurements. These primarily apply to benchmarks with a short execution timer of microseconds or less. These fluctuations can be caused by various factors, including environmental temperature, network traffic, \acrshort{os}-configuration, and others \parencite{chen2016_robust_benchmarking_noisy_environments}. In addition to these factors that directly relate to \acrshort{os} behaviour, variations arising from specific programming language features or implementation details are also present, such as errors due to timer inaccuracy \parencite{chen2016_robust_benchmarking_noisy_environments}.

One consequence of these variations is that the timing measurements are not necessarily \MakeLowercase{\acrfull{iid}}. Techniques such as F-tests and Student's t-tests have demonstrated low statistical power for non-\acrshort{iid} timing measurements. Parametric outlier detection techniques, such as the 3-sigma rule, may also fail when applied to non-\acrshort{iid} timing measurements \parencite{chen2016_robust_benchmarking_noisy_environments, halkos2018_corporate_social_responsibility}.

To ensure statistically reliable benchmarks, these variations must be substantially reduced. According to \textcite{chen2016_robust_benchmarking_noisy_environments}, variations can be reduced by increasing the benchmark run-time or repeating its execution multiple times. The authors argue that variations caused by factors outside the run-time environment become negligible as run-time increases, because the random delay becomes small relative to the total run-time. The timer inaccuracy is large relative to the single-run duration, but it is spread across all runs when multiple runs are measured. As a result, the timer inaccuracy is negligible when using the minimum estimator for the run-time. External variations can be neglected when the minimum estimator of the true run-time is used. \textcite{chen2016_robust_benchmarking_noisy_environments} state that variation never reduces run-time, and argue that the fastest run-time is also the one with the lowest variation.

\subsection{Variations and error sources in cloud environments}
\subsection{Measuring network traffic}
\subsection{Measuring cost in cloud environments}
\subsection{Sample size and statistical power}